\section{A general framework of finite interfaces and resources}
\label{sec:framework}

The method has two layers. A mathematical object is described by finite
observations and laws about them. A computation is justified by an
arithmetic program, a finite demand for the observations it uses, and a
bounded invariant. The program's termination resource is kept separate
from its mathematical correctness. This separation lets the same
representation support a predetermined recursion, an adaptive ranked
search, or an exhaustive certificate enumeration.

\subsection{Coding conventions}

Natural numbers are the only object-theory sort. Signed integers may be
coded by \(0\mapsto0\), \(n>0\mapsto2n\), and
\(-n\mapsto2n-1\). A rational is a pair consisting of an integer numerator
and a positive denominator; reduce by a bounded Euclidean algorithm when
a canonical code is needed. Invalid denominator codes are assigned a
fixed default value, while every mathematical theorem retains the
positive-denominator guard. Complex rational numbers are pairs of
rationals. The coding of finite lists supplies total primitive-recursive
length, lookup, append, and bounded fold functions. Values outside a
list's range have a default value and are never used without a range
test in a mathematical contract.

For definiteness, pairing can be coded by
\(\langle a,b\rangle=2^a(2b+1)-1\). Its projections use bounded search
below the code plus one. A list can be a length and a nested pairing,
with list traversal bounded by that length. Every concrete loop below
has a primitive-recursive bound; efficient binary encodings can replace
these without changing the arguments. Rational arithmetic, coefficient
arrays, determinant permutation sums, polynomial substitution, and
finite coefficient comparison are consequently primitive recursive.
There is no appeal here to a universal evaluator for arbitrary programs.

\begin{lemma}[Boolean compilation of fixed bounded formulas]
\label{lem:boolean}
For each fixed bounded formula \(A(\vec x)\) over primitive-recursive
symbols, there is a primitive-recursive \(a(\vec x)\in\bool\) such that
\(\ipra\) proves
\[
             A(\vec x)\ \longleftrightarrow\ a(\vec x)=1.
\]
The construction also supplies proofs of the range condition and of the
equivalence. It is a transformation on a fixed formula, not an internal
truth predicate for all arithmetic formulas.
\end{lemma}

\begin{proof}
Equality has a primitive-recursive comparison test, whose correctness
follows from the defining equations and induction on subtraction or
simultaneous predecessor comparisons. Atomic predicates are treated by
their fixed characteristic definitions. For Boolean values, conjunction
is multiplication, disjunction is the capped sum, and negation is
\(1-a\). Since the induction hypothesis proves the subformulas
decidable, implication is represented by
\(\max(1-a,b)\). This use of decidability is restricted to these
subformulas; no excluded-middle schema is introduced.

For \(\forall i<t(\vec x)\,B(i,\vec x)\), define the prefix product
\[
 u(0,\vec x)=1,\qquad
 u(n+1,\vec x)=u(n,\vec x)b(n,\vec x).
\]
Quantifier-free induction proves that the product remains Boolean.
To obtain a particular factor from \(u(n)=1\), fix its index \(i\)
as a parameter and induct on \(n\) with the quantifier-free formula
\(u(n)=1\to(i<n\to b(i)=1)\). In the successor case, a product
of two Boolean factors equals one only when both factors equal one;
compare the fixed index with the preceding prefix length.
For the converse, retain \(\forall i<t\,B(i)\) in the fixed context;
induct on \(n\) with \(n\le t\to u(n)=1\), using just the instance
at \(i=n\) in the successor case. This proves the required equivalence
without induction on the bounded universal formula itself.

For a bounded existential, use the capped prefix sum, or the
primitive-recursive first-success selector with a success bit. An
induction on the scan length maintains the Boolean equation saying
that the success bit agrees with the capped sum, and, when successful,
that the selected index is below the scan length and satisfies the
characteristic test. The forward existential implication instantiates
the product or sum at a supplied witness; the reverse implication uses
the selected index. The formula induction is a finite metatheoretic
construction, and all numerical inductions it emits are quantifier-free.
\end{proof}

\begin{corollary}[Bounded induction and bounded choice]
\label{cor:bounded}
Induction on any fixed bounded formula is derivable in \(\ipra\).
If \(\exists y\le B(\vec x)\,R(\vec x,y)\) is proved with bounded
\(R\), the bounded first-success scan supplies a fixed
primitive-recursive witness with an internal correctness proof.
Finite products of such scans give finite choice over bounded domains.
\end{corollary}

\begin{proof}
Replace the induction formula by the equation in
Lemma~\ref{lem:boolean}, perform quantifier-free induction, and replace
back. For choice, use the scan invariant just proved; the supplied
existence theorem forces its success bit to be one. Default outputs
make the function total even when the premise is false.
\end{proof}

\subsection{Finite observations}

\begin{definition}[An arithmetic observation interface]
Fix a primitive-recursive producer \(O(\vec a,k)\) of finite codes,
a bounded validity predicate \(V(\vec a,k,z)\), and bounded relations
for comparing or restricting observations. The parameters \(\vec a\)
describe an input presentation; \(k\) is an observation demand, not
necessarily a decimal or binary precision. An object assertion is an
ordinary first-order formula built from these contracts, for example
\[
   \forall k\,V(\vec a,k,O(\vec a,k)).
\]
No object of an infinite set or function space is thereby added to the
arithmetic domain.
\end{definition}

An observation may be a jet through degree \(k\), a matrix entry table,
a rational approximation with error \(2^{-k}\), a finite graph, a finite
list of roots without labels, or a polynomial identity certificate.
An equality of represented objects can be a family of bounded
observational equalities. An apartness assertion must carry an explicit
positive lower bound when an inverse is computed.

For instance, a real presentation can use a fixed rational-valued
program \(A(\vec a,k)\) and the regularity contract
\[
 |A(\vec a,p)-A(\vec a,q)|\le 2^{-p}+2^{-q}.
\]
Each instance is an integer inequality after clearing positive
denominators. One may reason under its universal closure, but cannot use
that universal closure as a quantifier-free induction formula. For a
fixed requested accuracy, the proof records exactly the finitely many
instances required. An equality test between arbitrary such names is
not a Boolean operation; approximate comparison with a specified margin
or a supplied exact algebraic code is a different interface.

\subsection{Operations and proof-bearing contracts}

An operation consists of fixed PR syntax for its numerical output,
fixed PR syntax for its finite input demands, and a derivation of its
local contract. For one input a typical contract is
\begin{equation}\label{eq:local-operation}
 V_A(\vec a,D(\vec a,k),z)
 \ \longrightarrow\
 V_B(\vec a,k,T(\vec a,k,z)).
\end{equation}
If several predecessor observations are needed, use a bounded list of
queries and the bounded conjunction of their contracts. If the producer
needs one demand for computing its output and a larger one for checking
the output, record both. Neither demand is inferred from the fact that
an ambient mathematical operation exists.

\begin{proposition}[Proof-producing composition]
\label{prop:composition}
Two operations of the form \eqref{eq:local-operation} compose. The
composite demand is obtained by substituting the second demand into the
first; the composite output is ordinary PR composition. Substitution
and implication elimination compose the two supplied internal proofs.
For a finite family of queries, a prefix maximum covers all demanded
indices, provided the validity interface explicitly supports the
corresponding prefix restriction.
\end{proposition}

\begin{proof}
If \(D_1\) supplies the input needed to produce an intermediate
observation at demand \(r\), and \(D_2\) is the demand for the final
request \(k\), use the first contract at \(r=D_2(k)\). Its conclusion
is exactly the premise of the second contract. PR syntax is closed
under the displayed substitutions. These are operations on given proof
terms; they do not use the semantic truth of an arbitrary input
certificate. For multiple queries the finite maximum dominates each
listed index by its defining recursion. Apply the restriction proof
before the local contract; absent that proof, a maximum is not a valid
replacement for several nonmonotone demands.
\end{proof}

\subsection{Execution resources and bounded invariants}

\begin{lemma}[Guarded finite execution]
\label{lem:execution}
Let \(u(0)=u_0\) and \(u(t+1)=T(\vec a,t,u(t))\) be an ordinary
primitive recursion, and let \(B(\vec a)\) be PR. Suppose a fixed context
\(\Gamma\), independent of the induction variable, proves a bounded
invariant \(I(\vec a,0,u_0)\) and the bounded-range step
\[
 t<B(\vec a),\ I(\vec a,t,u(t))
       \ \longrightarrow\ I(\vec a,t+1,u(t+1)).
\]
Then \(\Gamma\) proves \(I(\vec a,B(\vec a),u(B(\vec a)))\).
The context itself need not consist of bounded formulas.
\end{lemma}

\begin{proof}
Booleanize \(I\), keeping \(\vec a\) fixed. Induct on
\(t\le B(\vec a)\to I(\vec a,t,u(t))\), expressed by the corresponding
quantifier-free Boolean implication. From \(t+1\le B\) obtain \(t<B\)
and the preceding invariant; outside that range the implication has no
obligation. Instantiate the resulting universal formula at \(t=B\).
\end{proof}

There are three ways to supply \(B\) in this report. A bounded search
comes with its enumeration bound. A predetermined observation recursion
comes with its requested terminal stage, while its changing observation
parameters are compiled backwards. An adaptive computation with a
fixed-dimensional lexicographic rank receives a bound from
Section~\ref{sec:descent}. This third case permits output-dependent
control flow, provided the transition itself is fixed PR and the local
rank and invariant proofs are supplied.

\begin{theorem}[Restricted compilation principle]
\label{thm:compilation}
Consider a finite mathematical development whose representations and
operations have the above interfaces. Allow PR composition, bounded
folds and scans, the predetermined demand recursions of
Section~\ref{sec:demand}, and the fixed-depth ranked executions of
Section~\ref{sec:descent}. Require internal local proofs for each
primitive contract, proofs of initialization and preservation of a bounded
invariant for every numerical loop, and a proof that its terminal
invariant implies the requested output contract.
Then every output contract obtained by these composition, induction, and
output rules has an \(\ipra\) derivation under the supplied input
hypotheses. Every output witness
has fixed PR syntax. The transformation is a metatheoretic construction
from this finite development and its local derivations.
\end{theorem}

\begin{proof}
Proceed through the finite construction. Composition uses
Proposition~\ref{prop:composition}; scans and folds use
Corollary~\ref{cor:bounded} and Lemma~\ref{lem:execution}; a changing
parameter is replaced by the explicit two-recursion program of
Section~\ref{sec:demand}; a ranked loop is replaced by a fixed-fuel
simulation using Section~\ref{sec:descent}. Every generated invariant
is bounded, and Lemma~\ref{lem:boolean} supplies its quantifier-free
form. Existential introduction uses the explicitly generated witness
term. Universal introduction over numerical parameters is performed
only after their finite proof has been constructed and only when those
parameters do not occur freely in the retained context (otherwise the
relevant assumptions must first be discharged). No other induction
principle, soundness reflection, or uniform evaluator is used.
\end{proof}

The principle is deliberately restricted. It does not say that every
true specification of a PR program is provable in \(\ipra\), nor that
every Lean program is PR, nor that all constructive mathematics falls
inside arithmetic. Those claims would conceal exactly the obligations
the framework is meant to expose.

\subsection{A useful test for a proposed embedding}

The following four questions are a compact mathematical review tool.
What finite code is returned? Which observations of the input justify
it? What fixed PR construction bounds its execution? Which bounded
invariant proves its contract? A fifth question is needed for a
completeness theorem: why must every solution be represented by a code
the finite procedure examines? Sections~\ref{sec:elimination} and
\ref{sec:applications} answer the fifth question using elimination
polynomials and bounded auxiliary equations. Mere cardinality does not
answer it.
